% ============================================================================
% foreword.tex — Foreword by Steve Wilhite
% ============================================================================
\forewordpage[images/SEAS_short_white.png]
  {images/Steve.jpg}
  {Foreword}
  {Steve Wilhite}
  {Executive Vice President, Energy \& Sustainability Practice}

\noindent{\displayfont\fontsize{24}{14}\selectfont\color{SEGreen}What is the climate price tag on the world's compute infrastructure?\par}
\vspace{6mm}

\begin{multicols}{2}\raggedcolumns
\color{SEDarkGreen}
\noindent The data center industry stands at an inflection point. Over the past three decades, we have built the physical layer of the digital economy---facility by facility, market by market---into an infrastructure class that now underpins every sector of the global economy. The scale of what comes next dwarfs what came before: the AI era demands not incremental expansion but a structural transformation of how we design, power, and operate these critical assets.

\vspace{4mm}

\noindent At Schneider Electric, we work alongside operators, investors, and grid planners who are making capital allocation decisions today that will determine the resilience of digital infrastructure through 2050. The question they face is not whether climate risk affects data centers---every operator who has managed through a winter storm, a heatwave, or a grid emergency knows the answer. The question is how much, where, and what to do about it.

\vspace{4mm}

\noindent This paper provides the first quantitative answer. By translating physical and transition climate hazards into enterprise value loss at facility level---across the global installed base and the prospective AI pipeline---it makes visible a dimension of risk that capital markets have not yet priced. The supply chain analysis is particularly consequential: the finding that indirect business interruption exceeds direct exposure by an order of magnitude reframes how we think about facility resilience.

\columnbreak

\noindent What gives this work practical weight is that it does not stop at diagnosis. The adaptation analysis demonstrates that structured investment reduces climate exposure by a third, with positive returns under every scenario tested. For an industry deploying trillions in long-lived assets, that is an actionable result.

\vspace{4mm}

\noindent The analytical framework developed here---and the Schneider Electric Research Institute's commitment to open, reproducible methodology---provides the foundation for integrating climate risk into the investment decisions that will shape our industry for the next quarter century.
\end{multicols}